\documentclass[11pt,letterpaper]{letter}
\usepackage[margin=1in]{geometry}
\usepackage[utf8]{inputenc}
\usepackage[T1]{fontenc}
\usepackage{url}

\signature{Ihor Kendiukhov\\
Department of Computer Science\\
University of Tuebingen\\
Sand~14, 72076 Tuebingen, Germany\\
kenduhov.ig@gmail.com}

\address{Ihor Kendiukhov\\
Department of Computer Science\\
University of Tuebingen\\
Sand~14, 72076 Tuebingen, Germany}

\begin{document}

\begin{letter}{Editorial Board\\
PLOS Computational Biology}

\opening{Dear Editors,}

I am pleased to submit the manuscript entitled ``Causal intervention
validation of gene regulatory signals in scGPT'' for consideration as a
Research Article in \textit{PLOS Computational Biology}.

Single-cell foundation models such as scGPT are increasingly used for
cell-type annotation, perturbation prediction, and drug target
prioritization, yet whether their internal representations encode genuine
gene regulatory structure remains an open question. This manuscript
introduces a causal intervention pipeline that directly tests this by
ablating and swapping transcription factor token values within scGPT and
measuring downstream effects on known target genes. We evaluate across
three tissues from Tabula Sapiens, an external lung atlas, and
CRISPR perturbation data, comparing causal scores against attention and
coexpression baselines, and use activation patching to localize
responsible model components.

The key findings include: (i) lung tissue shows strong enrichment of
causal intervention scores for curated TF--target edges (AUPR up to 0.78,
permutation $p=0.001$), substantially exceeding attention and coexpression
baselines; (ii) kidney enrichment is modest but reproducible across
seeds, while immune tissue remains at baseline, revealing tissue-dependent
regulatory encoding; (iii) mechanistic tracing localizes effects to a
small set of attention heads and MLP blocks with depth separation between
self-targeting and cross-gene interactions; and (iv) external perturbation
validation provides an independent check, though with limited statistical
power at current scale.

I believe this work is well-suited for \textit{PLOS Computational Biology}
because it addresses the intersection of computational methodology
(causal inference and mechanistic interpretability applied to biological
foundation models) and biological insight (tissue-specific regulatory
encoding). The manuscript presents both the pipeline and an honest
assessment of where the approach succeeds and where it does not, which I
believe serves the community better than reporting only the strongest
results.

Source code and evaluation artifacts will be made publicly available upon
acceptance at \url{https://github.com/kendiukhov/biodyn}.

This manuscript has not been published elsewhere and is not under
consideration at any other journal. No conflicts of interest exist. I
declare that large language models were used for writing assistance and
code generation during the preparation of this manuscript.

Thank you for considering this submission.

\closing{Sincerely,}

\end{letter}
\end{document}
